\documentclass[11pt]{article}

\usepackage[margin=1in]{geometry}
\usepackage{amsmath,amssymb}
\usepackage{booktabs}
\usepackage{siunitx}
\usepackage[round,authoryear]{natbib}
\usepackage[colorlinks=true,allcolors=blue]{hyperref}

\newcommand{\Mnet}{\mathcal{M}_\theta}
\newcommand{\grid}{\ell}
\newcommand{\dhat}{\hat{\delta}}
\newcommand{\rhohat}{\hat{\rho}}
\newcommand{\shat}{\hat{\sigma}}
\newcommand{\Ahat}{\hat{A}}
\newcommand{\Nd}{N_d}
\newcommand{\Npix}{N_{\mathrm{pix}}}

\title{A Mat\'ern-class regularization term for training against sparse observations}
\date{\today}

\begin{document}
\maketitle

\begin{abstract}
\noindent
The network diagnoses a dense sea level anomaly field but is scored only against sparse
along-track observations, so at any cell no mission observed in the target bin the loss
carries no gradient at all. This note adds the missing term: the background term of a
variational data assimilation cost function, built from the Mat\'ern-class correlation
operator of \citet{smith2022}, so that the unobserved cells are filled by a controllable
correlation structure with a probabilistic interpretation rather than by whatever the
network extrapolates. Notation follows \citet{smith2022} throughout. The implementation
notes in Section~\ref{sec:impl} record three things the discretization on a land-masked
HEALPix mesh turns up: $\Phi$ is not isotropic, the operator retains a residual
anisotropy that we do not correct, and the stencil is first-order inconsistent at
coastal cells.
\end{abstract}

\section{The training loss}

\begin{equation}
\label{eq:loss}
\mathcal{L} =
\tfrac{1}{2}\left\| H\!\left(\Mnet(x_t, x_{t-\Delta t}, \ldots)\right) - y_{t+\Delta t}\right\|_2^2
\;+\; \frac{\lambda_w}{\alpha\,2}\left\|\theta\right\|_2^2
\;+\; \underbrace{\frac{\lambda_M}{2}\left\| C^{-1/2}\,\Mnet(x_t, x_{t-\Delta t}, \ldots)\right\|_2^2}_{\displaystyle \mathcal{L}_M}
\end{equation}

\noindent with

\begin{center}
\begin{tabular}{ll}
\toprule
$\theta$ & parameters \\
$\Mnet$ & neural network \\
$x_t$ & input state \\
$y_{t+\Delta t}$ & observations \\
$H$ & interpolation (observation operator) \\
$\lambda_w$ & regularization strength for weights \\
$\alpha$ & learning rate \\
$\lambda_M$ & regularization parameter, hopefully not needed \\
\bottomrule
\end{tabular}
\end{center}

\begin{description}
  \item[Weight decay.] The second term of \eqref{eq:loss} corresponds to weight decay
  only under plain SGD. We use AdamW, so it is written here for illustration and is not
  literally the term being minimised.

  \item[Background state.] $\mathcal{L}_M$ formally carries $(x-\,x_b)$ rather
      than $x$ inside the norm. We
  set $x_b := 0$, which in normalised space is the mean state. This is the only available
  choice: the predicted field is diagnostic, and the inputs are the sparse per-mission
  columns, so no dense background field exists to subtract.
\end{description}

\paragraph{Notation.} Two collisions with \citet{smith2022} are worth naming, since this
note keeps that paper's symbols everywhere else. First, $\theta(s)$ is the random field
in \citet{smith2022} but the parameters here, so the field is written $x$. Second, $M$ is
the exponent of the elliptic operator, and the network is $\Mnet$; the subscript
distinguishes them. Third, $H$ is the observation operator, so the grid scale introduced
in Section~\ref{sec:phi} is written $\grid$ rather than $h$.

\section{The correlation operator}

Following \citet{smith2022}, the correlation matrix is specified through its square root,

\begin{align}
C^{1/2}  &:= \Lambda A^{-M}\det(\Phi)^{-1/2}, \label{eq:chalf}\\
C^{-1/2} &\phantom{:}= \det(\Phi)^{1/2} A^{M} \Lambda^{-1}
                     = \det(\Phi)^{1/2} A^{M} \Sigma,
\qquad \Sigma := \Lambda^{-1}, \label{eq:cinvhalf}
\end{align}

\noindent where the elliptic operator and its parameters are

\begin{equation}
\label{eq:A}
A = \delta(s) - \nabla\cdot K(s) \nabla,
\qquad
K = \frac{\Phi\Phi^{T}}{\det(\Phi)},
\qquad
\delta(s) = \frac{\dhat}{\det(\Phi)},
\end{equation}

\begin{equation}
\label{eq:calib}
\dhat = \frac{8\nu}{\rhohat^{2}},
\qquad
\nu = 2M - \frac{\Nd}{2}.
\end{equation}

For a surface field $\Nd = 2$, so $\nu = 2M-1$: $M=1$ gives $\nu=1$ and $M=2$ gives
$\nu=3$. The half-integer case $M=\tfrac12$ is not available in two dimensions, since it
gives $\nu=0$. The range parameter $\rhohat$ is the distance at which correlation falls to
$\approx 0.14$, measured in grid cells.

Only \emph{forward} applications of $A$ are needed. This is the practical reason the term
is affordable: sampling from the prior would require $A^{-M}$ and hence a linear solve,
but evaluating $\mathcal{L}_M$ requires only $A^{M}$, which is a sparse matrix--vector
product.

\subsection{The normalization matrix}
\label{sec:lambda}

$\Lambda$ is fixed by requiring $C$ to be a correlation matrix, i.e.\ to have unit
diagonal. It is a diagonal matrix throughout, in the same form as \citet{smith2022},

\begin{equation}
\label{eq:lambda}
\Lambda = \det(\Phi)^{1/2}\operatorname{diag}\left\{1/\shat_n\right\}_{n=1}^{N},
\qquad
\Sigma = \Lambda^{-1} = \det(\Phi)^{-1/2}\operatorname{diag}\left\{\shat_n\right\}_{n=1}^{N},
\end{equation}

\noindent where $\shat_n$ is the pointwise standard deviation of $A^{-M}$ applied to unit
white noise, $\shat_n^2 \equiv (A^{-2M})_{nn}$. Substituting into \eqref{eq:cinvhalf} gives
$C_{nn} = \det(\Phi)^{-1}\Lambda_{nn}^2\shat_n^2 = 1$ as required, and the
$\det(\Phi)^{1/2}$ in $C^{-1/2}$ cancels against the one inside $\Sigma$:

\begin{equation}
\label{eq:LM}
\mathcal{L}_M = \frac{\lambda_M}{2}\left\|A^{M}\Sigma\, x\right\|_2^2
\qquad\text{with}\qquad
\Sigma \, x = \left(\shat_1 x_1, \; \shat_2 x_2, \; \ldots, \; \shat_N x_N\right)^{T}.
\end{equation}

Note $\Sigma$ acts on $x$ \emph{before} the operator is applied, not after. This is worth
stating because the alternative reading, $\|A^M x / \shat\|^2$, is a different quantity
and fails the checks below.

\paragraph{$\shat_n$ is a field here, not a scalar.} In the stationary, isotropic,
unbounded case the marginal variance has the closed form of \citet[Eq.~8]{smith2022},

\begin{equation}
\label{eq:analytic}
\shat^2 = \frac{\Gamma(\nu)}{\Gamma\!\left(\nu + \Nd/2\right)\left(4\pi\right)^{\Nd/2}\dhat^{\,\nu}},
\end{equation}

which depends only on $\dhat$, $\nu$ and $\Nd$ --- a single number, carrying no geometry.
That is the right object on an unbounded equal-area mesh, and Table~\ref{tab:analytic}
confirms it: run on a full sphere with no land mask, the estimated $\shat_n$ is spatially
constant to within the sampling noise of the estimator and agrees with \eqref{eq:analytic}
to between $1.4\%$ and $12\%$, converging as $\rhohat$ grows and the operator is better
resolved.

What breaks the scalar is the land mask. Dropping fluxes at a coastal face weakens $A$
locally, so $\shat_n$ acquires spatial structure that \eqref{eq:analytic} cannot represent
--- by a factor of up to $20$, as Section~\ref{sec:lowneighbour} shows. We therefore
estimate $\shat_n$ per cell by randomization \citep{weaver2021},
$\shat^2_n = \mathbb{E}[(A^{-M}z)_n^2]$ with $z \sim \mathcal{N}(0,I)$, which costs a
single sparse factorization plus $10^3$ solves, about \SI{3.4}{\second}, once at
initialization. Where the analytic scalar is adequate it is recovered automatically, so
nothing is lost by estimating.

\begin{table}[htbp]
\centering
\caption{Estimated $\shat_n$ against the analytic scalar \eqref{eq:analytic} on a full
sphere at $\mathrm{nside}=32$, no land mask. ``Spread'' is the $5$--$95$ percentile range
of $\shat_n$ over the mesh as a fraction of its median; with $600$ randomization samples
the estimator's own noise accounts for $\approx 9.5\%$, so $\shat_n$ is constant to within
the estimate. Reproduced by \texttt{scripts/check\_analytic.py}.}
\label{tab:analytic}
\begin{tabular}{cccccc}
\toprule
$M$ & $\rhohat$ & $\dhat$ & estimated $\shat_n$ & spread & estimated / analytic \\
\midrule
1 & 4  & 0.5000 & 0.44743 & 10.2\% & 1.122 \\
1 & 8  & 0.1250 & 0.84350 & 10.5\% & 1.057 \\
1 & 16 & 0.0312 & 1.63995 & 10.7\% & 1.028 \\
2 & 4  & 1.5000 & 0.09895 & 10.4\% & 1.116 \\
2 & 8  & 0.3750 & 0.72939 & 10.7\% & 1.028 \\
2 & 16 & 0.0938 & 5.75555 & 10.7\% & 1.014 \\
\bottomrule
\end{tabular}
\end{table}

\paragraph{Verification.} With $\Sigma$ so defined, drawing $x \sim \mathcal{N}(0,C)$ on
the real land-masked mesh gives marginal variance $0.998$ at interior cells and $0.986$ at
isolated cells against a target of $1$, and
$\mathbb{E}\!\left[x^{T} C^{-1} x\right] = 33436$ against $N = 33469$. Per-cell, the
contributions $\left(\left(A^{M}\Sigma x\right)_n\right)^2$ average to $0.999$, $0.999$,
$1.001$ and $0.972$ across cells with $8$, $3$--$7$, $1$--$2$ and $0$ retained neighbours
respectively.

\paragraph{What the code implements.} Two departures from \eqref{eq:loss}, both
deliberate. The implementation omits the $1/2$, because the paired data term
(\texttt{ObservationMaskedMSE}) is a plain MSE without one and carrying the factor on
only one term would be inconsistent. And $\lambda_M$ is not a parameter of the
regularization term: it is \texttt{loss\_weights[1]} of the \texttt{CombinedLoss} that
holds the two terms, so the strength lives in one place rather than two. What the code
evaluates is therefore

\begin{equation}
\label{eq:implemented}
\texttt{loss\_weights[1]} \times \frac{1}{N}\left\|\Ahat^{M}\hat{\Sigma}\,x\right\|_2^2 ,
\end{equation}

and anyone wanting the $1/2$ folds it into the weight.

\paragraph{Scaling.} Because \eqref{eq:LM} has expectation $N$ under the prior, defining
the implemented term as $\tfrac{\lambda_M}{2}\,\|\cdot\|^2 / N$ makes
$\mathbb{E}[\mathcal{L}_M] = \lambda_M/2$ at any mesh resolution and any $\rhohat$. Paired
with specifying the range in physical units and converting via
$\rhohat = \rho/(R\,\grid)$, this makes $\lambda_M$ dimensionless and transferable
between resolutions.

\section{Implementation notes}
\label{sec:impl}

The mesh is HEALPix level 6 ($\mathrm{nside}=64$, $\Npix = 49152$) in \textsc{nested}
ordering, pruned to the $N = 33469$ cells with \texttt{land\_fraction} $< 0.1$. The
discrete operator uses the eight-neighbour stencil from
\texttt{healpy.get\_all\_neighbours}, with land neighbours dropped --- the zero-flux
Neumann condition of \citet{smith2022}.

\paragraph{Index convention.} Throughout this section $i$ and $j$ index \emph{grid
cells}, never variables or spatial dimensions. Specifically $i$ is the cell whose
operator row is being written, and $j$ runs over the neighbours of $i$ that survived the
land mask --- at most the eight returned by \texttt{get\_all\_neighbours}, and fewer at
the coast. So $\vec{d}_{ij}$ is the offset from the centre of cell $i$ to the centre of
cell $j$, $\hat{d}_{ij} = |\vec{d}_{ij}|/\grid$ is its length in units of the grid
scale, $s_{ij}$ is the length of the face they share, and $w_{ij}$ is the weight on that
pair. Every $\sum_j$ below is over that short neighbour list, not over the whole mesh,
which is what makes the operator sparse.

\subsection{Choosing $\Phi$}
\label{sec:phi}

We adopt the same simplification as \citet[\S3.2]{smith2022}: take $\Phi$ diagonal,
ignoring its off-diagonal elements, with the normalizing length scales
$L_x(s)$, $L_y(s)$ on the diagonal. In two dimensions,

\begin{equation}
\label{eq:Phi}
\Phi = \begin{pmatrix} L_x & 0 \\ 0 & L_y \end{pmatrix},
\qquad
\det\Phi = L_x L_y,
\qquad
K = \frac{\Phi\Phi^{T}}{\det\Phi}
  = \begin{pmatrix} L_x/L_y & 0 \\ 0 & L_y/L_x \end{pmatrix},
\end{equation}

\noindent so $K$ carries the local aspect ratio and its reciprocal, and $\det K = 1$.
With $\delta(s) = \dhat/\det\Phi$ from \eqref{eq:A}, the operator is

\begin{equation}
\label{eq:A2d}
A = \frac{\dhat}{L_x L_y}
  - \left(\frac{L_x}{L_y}\frac{\partial^2}{\partial x^2}
        + \frac{L_y}{L_x}\frac{\partial^2}{\partial y^2}\right).
\end{equation}

This is the balance \citet{smith2022} is after: if the field varies on scale $L_x$ in
$x$ and $L_y$ in $y$, then both second-derivative terms are $O(x/(L_x L_y))$, the same
order as the zeroth-order term, so no direction dominates.

\paragraph{The nondimensional operator.} Multiplying \eqref{eq:A2d} through by
$\det\Phi$ removes the dimensional prefactor entirely:

\begin{equation}
\label{eq:Ahat}
\Ahat \;\equiv\; \det(\Phi)\,A
  \;=\; \dhat - \left(L_x^2 \frac{\partial^2}{\partial x^2}
                    + L_y^2 \frac{\partial^2}{\partial y^2}\right).
\end{equation}

$\Ahat$ is what the code actually builds. This is worth being explicit about,
because it is where the grid scale enters: the length scales appear as the coefficients
$L_x^2, L_y^2$ multiplying the second derivatives, and nowhere else.

The hat follows \citet{smith2022}: hatted quantities live in the grid-scale normalized
space, where the operator is isotropic and stationary. $\dhat$, $\rhohat$ and $\shat_n$
are already in that space, and $\Ahat$ and $\hat{d}_{ij}$ join them.

Nothing is lost by working with $\Ahat$ instead of $A$. Since
$A^{-M} = \det(\Phi)^{M}\Ahat^{-M}$, the square root \eqref{eq:chalf} collapses to
$C^{1/2} = \hat{\Lambda}\,\Ahat^{-M}$ with
$\hat{\Lambda} \equiv \det(\Phi)^{M-1/2}\Lambda$, and the unit-diagonal condition of
Section~\ref{sec:lambda} becomes simply

\begin{equation}
\label{eq:lambdahat}
\hat{\Lambda} = \operatorname{diag}\left\{1/\shat_n\right\}_{n=1}^{N},
\qquad
\hat{\Sigma} = \hat{\Lambda}^{-1} = \operatorname{diag}\left\{\shat_n\right\}_{n=1}^{N},
\qquad
\shat^2_n \equiv \left(\Ahat^{-2M}\right)_{nn}.
\end{equation}

Every power of $\det\Phi$ then cancels:

\begin{equation}
\label{eq:LMhat}
\mathcal{L}_M
 = \frac{\lambda_M}{2}\left\|\det(\Phi)^{1/2} A^{M} \Sigma x\right\|_2^2
 = \frac{\lambda_M}{2}\left\|\Ahat^{M}\hat{\Sigma}\, x\right\|_2^2 .
\end{equation}

So $\det\Phi$ never has to be evaluated at all, provided the normalization is estimated
for the same operator that is applied. This is the form the code implements: build
$\Ahat$, estimate $\shat_n$ from it, and apply \eqref{eq:LMhat}.

\subsection{The normalizing length scales on HEALPix}
\label{sec:lengths}

Following \citet[\S3.3]{smith2022} we take $L_x, L_y$ from the grid itself. HEALPix gives
one of these exactly and the other only approximately.

\paragraph{$\det\Phi$ is exact.} Every HEALPix pixel has area exactly $4\pi/\Npix$
(verified: \texttt{healpy.nside2pixarea} returns $2.55663465\times10^{-4}$~sr for every
pixel at $\mathrm{nside}=64$, identical to $4\pi/\Npix$). Since $\det\Phi = L_x L_y$ is
the pixel area, $\det\Phi = 4\pi/\Npix$ holds exactly, whatever the pixel shape.

\paragraph{The isotropic split is an approximation.} Splitting that area as
$L_x = L_y = \grid$ with

\begin{equation}
\label{eq:gridscale}
\grid \equiv \sqrt{4\pi/\Npix}
\end{equation}

gives $K = I$ by \eqref{eq:Phi}, which is what makes the implementation simple. But
$\grid$ is the square root of the pixel \emph{area}, not a pixel side length. HEALPix
pixels are equal-area curvilinear quadrilaterals whose shape varies with latitude, and
they are not squares. \texttt{healpy}'s own documentation for \texttt{nside2resol}, which
returns exactly this quantity, describes it as ``a gross approximation given the
different pixel shapes''.\footnote{\url{https://healpy.readthedocs.io/en/latest/generated/healpy.pixelfunc.nside2resol.html}}

Table~\ref{tab:sides} measures the true side lengths from
\texttt{healpy.boundaries}. In the equatorial zone ($|\mathrm{lat}| < 41.8^\circ$, the
four base pixels that tile the equatorial belt) the pixels are square to within
$0.2\%$ and their sides are within $4\%$ of $\grid$, so $L_x = L_y = \grid$ is a good
approximation there. In the polar caps the sides run from $0.80\,\grid$ to $1.46\,\grid$
with a within-pixel aspect ratio up to $1.48$, and the isotropic split is genuinely
approximate.

\begin{table}[htbp]
\centering
\caption{True HEALPix pixel side lengths at $\mathrm{nside}=64$, from
\texttt{healpy.boundaries}, in units of $\grid = \sqrt{4\pi/\Npix}$. The last column is
the ratio of longest to shortest side \emph{within} a pixel, which would be $1$ for a
square. Reproduced by \texttt{scripts/table\_pixel\_shape.py}.}
\label{tab:sides}
\begin{tabular}{lcccc}
\toprule
$|\mathrm{lat}|$ & min & median & max & within-pixel aspect \\
\midrule
$0$--$15$    & 1.002 & 1.005 & 1.007 & 1.000 \\
$15$--$30$   & 1.000 & 1.001 & 1.003 & 1.000 \\
$30$--$41.8$ & 1.003 & 1.015 & 1.042 & 1.002 \\
$41.8$--$60$ & 0.826 & 1.042 & 1.453 & 1.368 \\
$60$--$75$   & 0.805 & 1.023 & 1.458 & 1.425 \\
$75$--$90$   & 0.798 & 1.017 & 1.455 & 1.475 \\
\bottomrule
\end{tabular}
\end{table}

Numerically, $\grid = \SI{0.9161}{\degree} = \SI{101.87}{\kilo\metre}$ at
$\mathrm{nside}=64$ and $\SI{1.8323}{\degree} = \SI{203.74}{\kilo\metre}$ at
$\mathrm{nside}=32$.

\textbf{We take $L_x = L_y = \grid$ and hence $K = I$}, accepting the polar-cap error.
Section~\ref{sec:aniso} measures what that costs in the discrete operator.

\subsection{Discretizing the operator}
\label{sec:discrete}

Following the finite-volume approach, integrate the elliptic term of \eqref{eq:Ahat}
over the control volume $V_i$ of cell $i$ and apply the divergence theorem:

\begin{equation}
\label{eq:fv}
\int_{V_i} \nabla\cdot K\nabla x \; \mathrm{d}V
 = \oint_{\partial V_i} \left(K\nabla x\right)\cdot \hat{n} \; \mathrm{d}S
 \approx \sum_{j} k_{ij}\, s_{ij} \, \frac{x_j - x_i}{d_{ij}},
\end{equation}

where the sum runs over the faces cell $i$ shares with its neighbours, $s_{ij}$ is the
shared face length, $d_{ij}$ the centre-to-centre distance, and the gradient normal to
the face is approximated by the two-point difference. Dividing by $V_i$ and taking
$K = I$,

\begin{equation}
\label{eq:tpfa}
\left(\grid^2 \nabla^2 x\right)_i
 \approx \frac{\grid^2}{V_i}\sum_j \frac{s_{ij}}{d_{ij}}\left(x_j - x_i\right)
 = \sum_j \frac{s_{ij}}{d_{ij}}\left(x_j - x_i\right),
\end{equation}

\noindent because $V_i = \grid^2$ exactly on an equal-area mesh. \emph{This is where the
grid scale enters and immediately leaves again.} It appears in \eqref{eq:Ahat} as the
coefficient $\grid^2$, and cancels against the cell area in \eqref{eq:tpfa}, leaving a
purely geometric stencil. The discrete operator is therefore

\begin{equation}
\label{eq:Adiscrete}
\left(\Ahat{}x\right)_i
 = \left(\dhat + \sum_j w_{ij}\right) x_i - \sum_j w_{ij} x_j ,
\end{equation}

and the zero-flux Neumann condition is imposed simply by omitting the faces that abut
land from the sum --- no flux crosses them, so they contribute nothing.

\paragraph{Choice of weights.} Equation~\eqref{eq:tpfa} suggests $w_{ij} = s_{ij}/d_{ij}$,
the two-point flux approximation, over the four faces each pixel actually shares. That
scheme is only second-order consistent when the line joining two cell centres is
orthogonal to their shared face, which HEALPix satisfies near the equator but not towards
the poles. We therefore use a least-squares consistent stencil over all eight neighbours
instead. Taylor expanding \eqref{eq:Adiscrete} about cell $i$ with
$\vec{d}_{ij}$ the offset to neighbour $j$ in units of $\grid$,

\begin{equation}
\sum_j w_{ij}\left(x_j - x_i\right)
 = \underbrace{\left(\sum_j w_{ij}\vec{d}_{ij}\right)\cdot\nabla x}_{\text{must vanish}}
 + \frac{1}{2}\sum_j w_{ij}\, \vec{d}_{ij}^{\,T} \mathcal{H} \vec{d}_{ij}
 + O(\vec{d}^{\,3}),
\end{equation}

with $\mathcal{H}$ the Hessian. Matching $\grid^2\nabla^2 x$ requires two conditions:

\begin{equation}
\label{eq:consistency}
\sum_j w_{ij}\vec{d}_{ij} = 0
\qquad\text{and}\qquad
T_i \equiv \frac{1}{2}\sum_j w_{ij}\, \vec{d}_{ij}\vec{d}_{ij}^{\,T} = I .
\end{equation}

The first holds automatically wherever the neighbours are symmetrically placed. For the
second, take $w_{ij} = c/|\vec{d}_{ij}|^2$; then
$T_i = \tfrac{c}{2}\sum_j \hat{u}_{ij}\hat{u}_{ij}^{T}$ depends only on the neighbour
\emph{directions}, and its trace gives $\tfrac{c}{2} n_i = 2$, so $c = 4/n_i$ with
$n_i = 8$ the nominal neighbour count:

\begin{equation}
\label{eq:weights}
\boxed{\;w_{ij} = \frac{4}{8\,\hat{d}_{ij}^{\,2}}\;}
\qquad \hat{d}_{ij} = |\vec{d}_{ij}|/\grid .
\end{equation}

On a square lattice this assigns $0.5$ to edge-adjacent and $0.25$ to corner neighbours,
and $0.5\,\grid^2\nabla^2 + 0.25\cdot 2\grid^2\nabla^2 = \grid^2\nabla^2$ as required; the
measured medians on the real mesh are $0.496$ and $0.258$. Because $w \propto
|\vec{d}|^{-2}$ makes $T_i$ blind to distance, the pixel elongation of
Table~\ref{tab:sides} cancels out of the isotropic part of the operator entirely.

Table~\ref{tab:weights} confirms that this is the better of the two schemes on this mesh:
\eqref{eq:weights} satisfies the trace condition of \eqref{eq:consistency} to four
decimal places everywhere, while the two-point flux scheme is off by $1.3\%$ overall and
$9.2\%$ in the polar caps, and is also the more anisotropic of the two.

\begin{table}[htbp]
\centering
\caption{Consistency of the two candidate stencils at $\mathrm{nside}=64$ on the full
sphere, from \eqref{eq:consistency}. $\mathrm{tr}(T)/2 = 1$ and
$\lambda_{\max}/\lambda_{\min} = 1$ would be exact. Medians.
Reproduced by \texttt{scripts/check\_weights.py}.}
\label{tab:weights}
\begin{tabular}{lcccc}
\toprule
 & \multicolumn{2}{c}{$\mathrm{tr}(T)/2$} & \multicolumn{2}{c}{$\lambda_{\max}/\lambda_{\min}$} \\
\cmidrule(lr){2-3}\cmidrule(lr){4-5}
 & equatorial & polar & equatorial & polar \\
\midrule
two-point flux, $w = s_{ij}/d_{ij}$, 4 faces & 1.0086 & 1.0919 & 1.302 & 2.340 \\
implemented, \eqref{eq:weights}, 8 neighbours & 0.9999 & 0.9999 & 1.140 & 1.675 \\
\bottomrule
\end{tabular}
\end{table}

A useful by-product of the face matching in Table~\ref{tab:weights}: the four shared
faces map to the \texttt{SW}, \texttt{NW}, \texttt{NE} and \texttt{SE} slots of
\texttt{healpy.get\_all\_neighbours}, so HEALPix's true grid axes lie on the
\emph{diagonals} of that compass labelling. Table~\ref{tab:phi} shows the consequence in
the neighbour spacings: the diagonal slots sit at $1.005$--$1.042\,\grid$ while the
compass-aligned ones are the far neighbours at $1.18$--$1.79\,\grid$.

\begin{table}[htbp]
\centering
\caption{Neighbour spacings and stencil anisotropy at $\mathrm{nside}=64$, by latitude
band. Distances are medians in units of $\grid$; $L_x/L_y$ is the ratio of the
east--west to north--south far-neighbour spacing. Reproduced by
\texttt{scripts/table\_phi.py}.}
\label{tab:phi}
\begin{tabular}{lccccc}
\toprule
$|\mathrm{lat}|$ & shared face & far nbr (E--W) & far nbr (N--S) & $L_x/L_y$ & $\lambda_{\max}/\lambda_{\min}$ \\
\midrule
$0$--$15$    & 1.005 & 1.523 & 1.313 & 1.160 & 1.159 \\
$15$--$30$   & 1.001 & 1.423 & 1.406 & 1.012 & 1.055 \\
$30$--$41.8$ & 1.015 & 1.252 & 1.590 & 0.791 & 1.260 \\
$41.8$--$60$ & 1.042 & 1.179 & 1.783 & 0.656 & 1.705 \\
$60$--$75$   & 1.022 & 1.227 & 1.732 & 0.714 & 1.675 \\
$75$--$90$   & 1.017 & 1.248 & 1.692 & 0.755 & 1.618 \\
\bottomrule
\end{tabular}
\end{table}

\subsection{Anisotropy}
\label{sec:aniso}

Section~\ref{sec:discrete} showed that the trace condition of
\eqref{eq:consistency} is met exactly. What does not follow is the traceless part: the
choice $L_x = L_y = \grid$ of Section~\ref{sec:lengths} sets $K = I$, but the pixels it
describes are not isotropic, and the stencil inherits that. The last column of
Table~\ref{tab:phi} gives the eigenvalue ratio $\lambda_{\max}/\lambda_{\min}$ of $T_i$:
between $1.06$ and $1.70$ depending on latitude, worst in the band $42^\circ$--$60^\circ$,
with median $1.18$ over the mesh and a $95$th percentile of $2.26$. It tracks the
within-pixel aspect ratio of Table~\ref{tab:sides}, as it should --- both are small in
the equatorial zone and appear together in the polar caps.

\textbf{We take $K = I$ regardless, as a stated assumption.} The consequence is that
correlation ellipses are elongated by up to $\sqrt{1.70}\approx 1.3$ in the mid-latitudes
rather than being circular. This is consistent with, and the origin of, the residual
directional spread of $\approx 0.04$ in correlation units reported in
Section~\ref{sec:validation}.

The justification is proportionality: $\lambda_M$ and $\rho$ are tuned quantities, and a
$1.3{:}1$ distortion of the prior's correlation ellipses is small next to the uncertainty
in what the correlation length should be in the first place. The fix, should it ever
matter, is available and cheap: set $K_i = T_i^{-1}$ from the measured second-moment
tensor, folded into the edge weights at initialisation, at the cost of making the weights
direction-dependent rather than a scalar per edge.

\subsection{Low-neighbour points}
\label{sec:lowneighbour}

Of the $N = 33469$ retained cells, $30064$ keep all eight neighbours, $3295$ have three to
seven, $98$ have one or two, and $12$ have none at all. The isolated cells are the
Kattegat, the Gulf of Bothnia, the White Sea, one cell in the Philippines, and eight in
the Canadian Arctic Archipelago and Baffin region. Their own \texttt{land\_fraction} is
$0.004$--$0.093$, so they are genuine ocean, ringed by cells that failed the threshold.

\paragraph{Diagnosis.} Dropping land fluxes makes the stencil \emph{first-order
inconsistent} at boundary cells. A consistent Laplacian must annihilate a linear field,
i.e.\ $\sum_j w_{ij}\vec{d}_{ij} = 0$; the measured median of
$\left|\sum_j w_{ij}\vec{d}_{ij}\right|$ is $0.0071$ at interior cells but $0.7847$ at
cells with three to seven neighbours. Such a cell therefore responds to a linear field at
$O(|\nabla x|\grid)$ where the interior responds at $O(\nabla^2 x\,\grid^2)$. The
Laplacian is simultaneously under-weighted there: $\mathrm{tr}(T)/2$ is $0.75$ at
three-to-seven-neighbour cells and $0.25$ at one-to-two-neighbour cells, against $1$ in
the interior.

\paragraph{Consequences.} Both effects inflate $\shat$, and the inflation grows with the
correlation length exactly as Fig.~4 of \citet{smith2022} shows (Table~\ref{tab:sigma}).
The penalty consequence is Table~\ref{tab:coast}. A field drawn \emph{from the prior}
gives unit contribution everywhere, so the normalization is exact by its own measure; the
excess is specific to fields that violate the Neumann boundary condition, which any
realistic ocean field does.

\begin{table}[htbp]
\centering
\caption{$\shat/\shat_{\text{interior}}$ at $M=1$, by retained-neighbour count and
correlation range. Reproduces Fig.~4 of \citet{smith2022}.}
\label{tab:sigma}
\begin{tabular}{lcccc}
\toprule
$\rhohat$ (cells) & 2 & 4 & 8 & 16 \\
range (\si{\kilo\metre}) & 204 & 407 & 815 & 1630 \\
\midrule
3--7 neighbours & 1.17 & 1.31 & 1.43 & 1.50 \\
1--2 neighbours & 1.71 & 2.54 & 3.44 & 3.75 \\
isolated        & 2.18 & 4.44 & 9.36 & 19.0 \\
\bottomrule
\end{tabular}
\end{table}

\begin{table}[htbp]
\centering
\caption{Per-cell penalty contribution relative to the interior, $M=1$, $\rhohat=4$, for a
generic smooth field with non-zero normal gradient at the coast.}
\label{tab:coast}
\begin{tabular}{lccc}
\toprule
 & 3--7 nbrs & 1--2 nbrs & isolated \\
\midrule
normalised, $\Sigma$ from \eqref{eq:lambda} & $2.14\times$ & $7.19\times$ & $13.9\times$ \\
unnormalised, $\Sigma$ constant ($\Lambda = I$) & $0.86\times$ & $0.72\times$ & $0.70\times$ \\
\bottomrule
\end{tabular}
\end{table}

Note that the two rows of Table~\ref{tab:coast} err in opposite directions. Leaving
$\Lambda = I$ does not avoid the problem; it under-constrains the same cells instead.

\paragraph{Decision.} \textbf{We drop the regularization term at low-neighbour points},
controlled by a \texttt{min\_neighbours} threshold. The default of $3$ removes the $110$
cells with fewer than three neighbours --- $0.33\%$ of the mesh, and the $7.2\times$ and
$13.9\times$ rows of Table~\ref{tab:coast} --- while retaining the $3295$ cells at
$2.14\times$. Raising the threshold to $8$ removes all $3405$ boundary cells, about
$10\%$ of the mesh.

The reasoning is that these cells do not merit more engineering than that. They are an
artefact of where the \texttt{land\_fraction} threshold happens to fall and would change
or disappear if it moved, and the Laplacian at a cell with no neighbours is not a
meaningful object in the first place. If dropping them proves too blunt --- for instance
if it leaves visible artefacts at the coast --- the next step is to cap the amplitude
instead, clipping $\shat$ at a fixed multiple of its interior value, which bounds the
excess in Table~\ref{tab:coast} while keeping the interior exact.

\paragraph{Alternatives considered and rejected.} Two more principled fixes were tested
and are recorded here so they are not re-derived later.

\begin{enumerate}
  \item \emph{Gradient-consistency projection.} Project $w_i$ onto the constraint
  $\sum_j w_{ij}\vec{d}_{ij} = 0$ and rescale to restore $\mathrm{tr}(T)/2 = 1$. This
  works for three-to-seven-neighbour cells, driving the gradient response to zero and the
  trace to $1.00$. But it annihilates one-to-two-neighbour cells entirely
  ($\mathrm{tr}(T)/2 \to 0$, since two constraints exhaust two degrees of freedom) and
  introduces $2500$ negative weights with minimum $-320$, which would very likely cost
  the positive-definiteness that makes $C$ a valid covariance. A non-negativity-constrained
  formulation might recover it, but was not pursued.

  \item \emph{Trace renormalisation at the boundary.} Rescale boundary weights so
  $\mathrm{tr}(T)/2 = 1$ everywhere, keeping all weights positive. This improves the
  three-to-seven-neighbour cells from $2.14\times$ to $1.51\times$ and their $\shat$ ratio
  from $1.31$ to $1.12$, but does nothing for isolated cells. It also makes the first-order
  inconsistency \emph{worse}, since scaling the weights up amplifies it:
  $\left|\sum_j w_{ij}\vec{d}_{ij}\right|$ rises from $0.78$ to $1.06$ at
  three-to-seven-neighbour cells and from $0.55$ to $3.20$ at one-to-two-neighbour cells.
  And it is no longer the strict zero-flux Neumann condition of \citet{smith2022} --- it
  is closer to a mirror boundary condition.
\end{enumerate}

\section{Validation against the analytical correlation function}
\label{sec:validation}

Because $A$ is symmetric, $\mathrm{Cov} = A^{-M}A^{-M,T} = A^{-2M}$ exactly, so no
sampling is needed to check the discretization. Taking a single column of $A^{-2M}$ by
sparse factorization and normalising it by its own diagonal gives the discrete
correlation as a function of distance from that source cell, which is compared against

\begin{equation}
\label{eq:matern}
r(d;\rhohat,M) = \frac{1}{2^{\nu-1}\Gamma(\nu)} z^{\nu} K_\nu(z),
\qquad
z = \sqrt{8\nu}\,\frac{d}{\rhohat}.
\end{equation}

Table~\ref{tab:validation} reports this on a full sphere at $\mathrm{nside}=32$ --- no
land mask, so that stencil quality is isolated from the boundary effects of
Section~\ref{sec:lowneighbour} --- from three source cells at different latitudes, since
HEALPix pixel shape varies with latitude even though area does not.

\begin{table}[htbp]
\centering
\caption{Discrete versus analytical correlation, $\mathrm{nside}=32$, full sphere. Entries
are equator / mid-latitude / near-pole source cells. The effective range
$\rho_{\text{emp}}$ is the distance at which the discrete correlation crosses the
analytical value $r(\rhohat)$. Directional spread is the standard deviation of the
discrete correlation across all cells in the band $0.8\rhohat < d < 1.2\rhohat$, which
would be zero for a perfectly isotropic operator.}
\label{tab:validation}
\begin{tabular}{ccccccc}
\toprule
$M$ & $\nu$ & $\rhohat$ & $r(\rhohat)$ & $\max|\text{err}|$ & $\rho_{\text{emp}}/\rhohat$ & directional spread \\
\midrule
1 & 1 & 4 & 0.140 & 0.13 / 0.15 / 0.14 & 0.91 / 0.88 / 0.85 & 0.031 / 0.034 / 0.034 \\
1 & 1 & 8 & 0.140 & 0.07 / 0.12 / 0.10 & 0.93 / 0.81 / 1.01 & 0.034 / 0.037 / 0.040 \\
2 & 3 & 4 & 0.138 & 0.13 / 0.16 / 0.16 & 0.94 / 0.88 / 0.85 & 0.039 / 0.041 / 0.043 \\
2 & 3 & 8 & 0.138 & 0.05 / 0.14 / 0.09 & 0.98 / 0.89 / 1.01 & 0.045 / 0.047 / 0.053 \\
\bottomrule
\end{tabular}
\end{table}

Three readings of Table~\ref{tab:validation}:

\begin{itemize}
  \item The effective range is within about $15\%$ of the requested $\rhohat$, so
  $\rhohat$ is a calibrated parameter rather than a proxy for one.

  \item The analytical $r(\rhohat)$ comes out at $0.140$ for $\nu=1$ and $0.138$ for
  $\nu=3$, the property asserted by Fig.~1 of \citet{smith2022}. This is the check that
  confirms the calibration $\dhat = 8\nu/\rhohat^2$ in \eqref{eq:calib}; dropping the
  $\nu$ breaks it for $M \neq 1$.

  \item The maximum absolute error is dominated by separations below two grid cells,
  where no discrete stencil can reproduce the continuum Green's function --- the $\nu=1$
  Mat\'ern has a cusp at the origin. It falls as $\rhohat$ grows, as expected.

  \item The directional spread of $0.03$--$0.05$ is the residual anisotropy of
  Section~\ref{sec:aniso} expressed in correlation units.
\end{itemize}

\begin{thebibliography}{9}

\bibitem[Smith(2022)]{smith2022}
T.~A. Smith.
\newblock A practical formulation for an anisotropic and nonstationary Mat\'ern class
correlation operator.
\newblock \emph{ESS Open Archive}, 2022.
\newblock \url{https://doi.org/10.1002/essoar.10511974.1}.

\bibitem[Lindgren et al.(2011)]{lindgren2011}
F.~Lindgren, H.~Rue, and J.~Lindstr\"om.
\newblock An explicit link between Gaussian fields and Gaussian Markov random fields: the
stochastic partial differential equation approach.
\newblock \emph{Journal of the Royal Statistical Society B}, 73(4):423--498, 2011.

\bibitem[Weaver et al.(2021)]{weaver2021}
A.~T. Weaver, M.~Chrust, B.~M\'en\'etrier, and A.~Piacentini.
\newblock An evaluation of methods for normalizing diffusion-based covariance operators in
variational data assimilation.
\newblock \emph{Quarterly Journal of the Royal Meteorological Society}, 147(734):289--320,
2021.

\end{thebibliography}

\end{document}
